Committing tokens on a per-position basis removes the chance for later revision.
We simulate this consequence of per-position commitment by freezing each token as soon as it is accepted by the EB rule.
The main property we study is that the token can no longer change.
We find that later revision becomes increasingly useful when trading off performance for speed at higher entropy bounds, and that self-correction persists without renoising.

\paragraph{Setup}
We focus on the EB sampler, as it is the default sampler used in \citet{team2026diffusiongemma}.
Our goal is to study the loss of sampler flexibility under per-position commitment, and determine whether later revisions require explicit noise injection.
To do this, we measure the macro-average of LCB v6, Natural2Code, and GPQA Diamond.
We vary two sampler choices: whether accepted positions are frozen, and whether rejected positions are renoised. The model weights and uniform corruption process are unchanged in all four conditions.

\paragraph{Frozen EB Sampler}
The Frozen EB ablation simulates per-position commitment: once a position is accepted under the EB rule, its token can no longer change. We retain the denoising window and blockwise cache commitment to isolate the loss of later revision. This simulation does not reproduce the attention changes or computational savings of actually moving individual positions into the KV cache.
We refer to the unmodified Unified EB sampler as ``unconstrained'', as any position can be changed regardless of whether it was already accepted under the EB rule.

\paragraph{No Noise}
To test whether revisions require explicit noise injection, we additionally disable renoising of rejected positions.
In this condition, rejected positions retain their current tokens, while accepted predictions can still revise unfrozen positions. This isolates the role of explicit renoising without changing the underlying model.

\paragraph{Ablation}
Removing noise and freezing accepted tokens are two independent interventions. Thus we ablate over all four combinations and show the results in Figure \ref{fig:frozen-eb-frontier}.
For each possible combination, we do a sweep over various entropy bound values to examine how they are affected by more relaxed thresholds.
For the unconstrained variants, we also track the empirical frequencies of each position being altered during the denoising trajectory after the initial acceptance under the EB rule.
We show these position volatility metrics in Figure~\ref{fig:wasted} (bottom row).

\paragraph{Results}
We observe that for each adaptive variant, the unconstrained sampler achieves a superior Pareto frontier than the frozen sampler, and this gap increases with further relaxation of the thresholds.
At the loosest bound, the frozen variants fall to about 20\% macro Pass@1, while the unconstrained variants stay near 70\%.

Furthermore, the number of positions that change after initial acceptance increases with higher entropy thresholds, and is relatively unaffected by whether the sampler uses noise injection.
At the default bound of $0.1$, roughly 97\% to 99\% of accepted positions never change, so the default sampler behaves much like a frozen one; revisions become frequent only as the bound is relaxed.
The simulation shows the cost of removing later revision on a per-position basis: accepting a token once does not mean that fixing it will preserve performance at higher entropy bounds. Self-correction also persists without renoising, so explicit noise injection is not necessary for these revisions.

Per-position commitment therefore loses sampler flexibility as the entropy bound is relaxed. This motivates waiting for more than an individual position to be accepted before removing its chance for later revision.
